\documentclass{beamer}
\usetheme{Madrid}
\usecolortheme{default}
\usepackage[utf8]{inputenc}
\usepackage{amsmath}
\usepackage{amsfonts}
\usepackage{amssymb}
\usepackage{tikz}
\usepackage{xcolor}
\usepackage{hyperref}

\title{CRISC Exam - Domain 4: Technology \& Security}
\subtitle{Technology Risk Questions 66--105 with Explanations}
\author{Compiled from ISACA CRISC Review Manual 8th Edition}
\date{\today}

\begin{document}
	
	\begin{frame}
		\titlepage
	\end{frame}
	
	\begin{frame}{Outline}
		\tableofcontents
	\end{frame}
	
	% ================ SECTION 6: EMERGING TECHNOLOGIES ================
	
	\section{Emerging Technologies}
	
	\begin{frame}{Q66: Shadow IT Risk}
		\begin{block}{Question}
			Why is shadow IT a significant risk to an enterprise?
		\end{block}
		\begin{enumerate}[(A)]
			\item It reduces IT costs
			\item It puts the organization in a vulnerable position where its true risk profile is unknown even to senior executives
			\item It improves security controls
			\item It eliminates the need for IT governance
		\end{enumerate}
	\end{frame}
	
	\begin{frame}{Q66: Explanation}
		\begin{block}{Correct Answer: B}
			If the IT department is unwilling to review and integrate new technologies, it is possible that the business will find ways to implement them regardless, putting the organization in a vulnerable position where its true risk profile is unknown even to senior executives. This practice, sometimes called shadow IT, is surprisingly common.
		\end{block}
		\begin{block}{Option Analysis}
			\begin{itemize}
				\item \textbf{A - Incorrect:} Shadow IT often increases costs
				\item \textbf{B - Correct:} True risk profile becomes unknown
				\item \textbf{C - Incorrect:} Shadow IT bypasses security controls
				\item \textbf{D - Incorrect:} It increases the need for governance
			\end{itemize}
		\end{block}
	\end{frame}
	
	\begin{frame}{Q67: Risk Practitioner and Emerging Tech}
		\begin{block}{Question}
			What is the BEST approach for a risk practitioner when new technologies emerge?
		\end{block}
		\begin{enumerate}[(A)]
			\item Adopt an obstructionist posture and reject all new technology
			\item Lead the effort to demonstrate how new technologies can be safely incorporated into the enterprise
			\item Ignore new technologies until they are mandated
			\item Delegate all technology decisions to the IT department
		\end{enumerate}
	\end{frame}
	
	\begin{frame}{Q67: Explanation}
		\begin{block}{Correct Answer: B}
			The risk practitioner is considerably more likely to get the attention and support of senior management by leading the effort to demonstrate how new technologies can be safely incorporated into the enterprise than by adopting an obstructionist posture. The best chance of maintaining relevance is to proactively suggest how technology can be put to use safely rather than seeking to reject its adoption.
		\end{block}
		\begin{block}{Option Analysis}
			\begin{itemize}
				\item \textbf{A - Incorrect:} Obstructionist posture loses influence
				\item \textbf{B - Correct:} Lead safe incorporation efforts
				\item \textbf{C - Incorrect:} Ignoring technology increases shadow IT
				\item \textbf{D - Incorrect:} Risk management is an enterprise responsibility
			\end{itemize}
		\end{block}
	\end{frame}
	
	\begin{frame}{Q68: Omnipresent Connectivity and Behavior}
		\begin{block}{Question}
			How has omnipresent connectivity changed human behavior in ways that create risk?
		\end{block}
		\begin{enumerate}[(A)]
			\item People are less likely to share information
			\item People are more likely to share updates about their actions, locations, and experiences on social media as they happen
			\item People have stopped using mobile devices
			\item People are more careful about information privacy
		\end{enumerate}
	\end{frame}
	
	\begin{frame}{Q68: Explanation}
		\begin{block}{Correct Answer: B}
			People are more likely to share updates about their actions, locations, and experiences on social media as they happen, including livestream video feeds that may be taken without notice to those around them. Many users of powerful mobile devices demonstrate minimal awareness of or even interest in basic concepts of cybersecurity and information privacy.
		\end{block}
		\begin{block}{Option Analysis}
			\begin{itemize}
				\item \textbf{A - Incorrect:} People are more likely to share
				\item \textbf{B - Correct:} Increased sharing creates risk
				\item \textbf{C - Incorrect:} Mobile device use has increased
				\item \textbf{D - Incorrect:} Many users show minimal privacy awareness
			\end{itemize}
		\end{block}
	\end{frame}
	
	\begin{frame}{Q69: IoT Device Security Priority}
		\begin{block}{Question}
			Why do IoT devices tend to have weaker security than traditional IT systems?
		\end{block}
		\begin{enumerate}[(A)]
			\item Security features make things more expensive and slow them down, and functional tasks are often seen as part of physical plant concerns
			\item IoT devices are always more secure
			\item IoT devices are managed by the security team from the start
			\item IoT devices do not connect to networks
		\end{enumerate}
	\end{frame}
	
	\begin{frame}{Q69: Explanation}
		\begin{block}{Correct Answer: A}
			IoT devices tend to prioritize functionality over security. The basic tension between operating efficiency and security is hard to overcome; security features make things more expensive and slow them down. Functional tasks linked to IoT devices also tend to be seen as part of the physical plant or facilities area of concern, receiving relatively little scrutiny within system life cycles or enterprise architecture.
		\end{block}
		\begin{block}{Option Analysis}
			\begin{itemize}
				\item \textbf{A - Correct:} Cost and speed trade-offs, seen as facilities concern
				\item \textbf{B - Incorrect:} IoT devices often have weaker security
				\item \textbf{C - Incorrect:} IoT is often managed by facilities, not security
				\item \textbf{D - Incorrect:} IoT devices are network-connected
			\end{itemize}
		\end{block}
	\end{frame}
	
	\begin{frame}{Q70: IoT Guardrails Approach}
		\begin{block}{Question}
			What is a recommended guardrails approach for managing IoT risk?
		\end{block}
		\begin{enumerate}[(A)]
			\item Ban all IoT devices
			\item Place IoT traffic on a separate network and treat that environment as hostile cyberterrain
			\item Connect IoT devices directly to critical systems
			\item Ignore IoT devices since they are low risk
		\end{enumerate}
	\end{frame}
	
	\begin{frame}{Q70: Explanation}
		\begin{block}{Correct Answer: B}
			Where practical, it may be beneficial to place IoT traffic on a separate network and treat that environment as hostile cyberterrain. This sort of guardrails approach can allow IT operations and security staff to advocate for the benefits of IoT while maintaining acceptable levels of risk, rather than creating conditions under which IoT capabilities might be brought into the enterprise as a form of shadow IT.
		\end{block}
		\begin{block}{Option Analysis}
			\begin{itemize}
				\item \textbf{A - Incorrect:} Banning all IoT is often impractical
				\item \textbf{B - Correct:} Separate network as hostile cyberterrain
				\item \textbf{C - Incorrect:} Direct connection increases risk
				\item \textbf{D - Incorrect:} Ignoring IoT increases risk
			\end{itemize}
		\end{block}
	\end{frame}
	
	\begin{frame}{Q71: Decryption and Massive Computing Power}
		\begin{block}{Question}
			How has cloud computing changed the risk associated with encryption?
		\end{block}
		\begin{enumerate}[(A)]
			\item Encryption is no longer needed
			\item It is now feasible for individuals or groups with modest means to perform brute-force attacks using rapidly provisioned computing power
			\item Encryption has become unbreakable
			\item Cloud providers cannot decrypt any data
		\end{enumerate}
	\end{frame}
	
	\begin{frame}{Q71: Explanation}
		\begin{block}{Correct Answer: B}
			Until recently, it was infeasible that anyone outside a well-funded government agency would be able to assemble sufficient processing power to break encryption. The ability to rapidly provision computing power through cloud providers has eliminated this restriction, with pattern analysis and brute-force attacks now within the grasp of individuals or groups with modest means.
		\end{block}
		\begin{block}{Option Analysis}
			\begin{itemize}
				\item \textbf{A - Incorrect:} Encryption remains vital
				\item \textbf{B - Correct:} Cloud enables brute-force attacks
				\item \textbf{C - Incorrect:} No encryption is unbreakable forever
				\item \textbf{D - Incorrect:} Providers can access unencrypted data
			\end{itemize}
		\end{block}
	\end{frame}
	
	\begin{frame}{Q72: Quantum Computing and Encryption}
		\begin{block}{Question}
			What does the emergence of quantum computing mean for encrypted data?
		\end{block}
		\begin{enumerate}[(A)]
			\item All encrypted data is safe forever
			\item It is difficult to predict how long an encrypted file might protect data if it were stolen
			\item Encryption will never need to change
			\item Quantum computing has no impact on cryptography
		\end{enumerate}
	\end{frame}
	
	\begin{frame}{Q72: Explanation}
		\begin{block}{Correct Answer: B}
			For the risk practitioner, this means that it is difficult to predict how long an encrypted file might protect data if it were stolen. Encryption remains a vital part of information security, but it cannot be taken for granted that data is safe simply because it is encrypted. Quantum-resistant algorithms have already begun appearing, and maintaining awareness of these offerings can help reduce the risk.
		\end{block}
		\begin{block}{Option Analysis}
			\begin{itemize}
				\item \textbf{A - Incorrect:} No encryption is safe forever
				\item \textbf{B - Correct:} Prediction of protection duration is difficult
				\item \textbf{C - Incorrect:} Encryption must evolve
				\item \textbf{D - Incorrect:} Quantum computing impacts cryptography significantly
			\end{itemize}
		\end{block}
	\end{frame}
	
	\begin{frame}{Q73: Deepfake Countermeasures}
		\begin{block}{Question}
			Which combination of controls can help guard against deepfakes misdirecting operations?
		\end{block}
		\begin{enumerate}[(A)]
			\item Ignoring all verbal approvals
			\item Security awareness and administrative controls such as requiring written approval or using a dedicated callback number
			\item Relying solely on video identification
			\item Eliminating all financial transfers
		\end{enumerate}
	\end{frame}
	
	\begin{frame}{Q73: Explanation}
		\begin{block}{Correct Answer: B}
			Risk practitioners can guard against deepfakes misdirecting operations through a combination of security awareness and administrative controls. For instance, requiring that a transaction not only have the verbal approval of the CEO but also written approval from another executive, or using a dedicated callback number to obtain confirmation after receiving initial outreach, can help inoculate an organization against this kind of technical deception.
		\end{block}
		\begin{block}{Option Analysis}
			\begin{itemize}
				\item \textbf{A - Incorrect:} Verbal approvals are still needed for operations
				\item \textbf{B - Correct:} Awareness plus administrative controls
				\item \textbf{C - Incorrect:} Video can be deepfaked
				\item \textbf{D - Incorrect:} Eliminating transfers is impractical
			\end{itemize}
		\end{block}
	\end{frame}
	
	\begin{frame}{Q74: Deepfake Reputational Risk}
		\begin{block}{Question}
			Why is deepfake technology a significant concern for reputational risk?
		\end{block}
		\begin{enumerate}[(A)]
			\item It only affects internal communications
			\item It can create falsified video or audio recordings that impact public opinion, particularly for enterprises subject to public tensions
			\item It is easy to detect by the average person
			\item It cannot be used for political or commercial gain
		\end{enumerate}
	\end{frame}
	
	\begin{frame}{Q74: Explanation}
		\begin{block}{Correct Answer: B}
			No protective measures are sufficient to guard against impacts to public opinion that may arise from falsified video or audio recordings, particularly when dealing with enterprises or issues subject to substantial public tensions. Many people for whom it is easiest to gather the data points needed to produce deepfakes are in positions where reputational risk is significant. Until there is an easy way to identify deepfakes, deception for political or commercial gain will remain an emerging trend requiring a robust risk response.
		\end{block}
		\begin{block}{Option Analysis}
			\begin{itemize}
				\item \textbf{A - Incorrect:} Deepfakes affect public perception externally
				\item \textbf{B - Correct:} Public opinion impacts are significant
				\item \textbf{C - Incorrect:} Detection requires considerable expertise
				\item \textbf{D - Incorrect:} Deepfakes can be used for gain
			\end{itemize}
		\end{block}
	\end{frame}
	
	\begin{frame}{Q75: Big Data Privacy - Reidentification}
		\begin{block}{Question}
			What privacy risk occurs when data is aggregated and semianonymous information becomes identifiable?
		\end{block}
		\begin{enumerate}[(A)]
			\item Data minimization
			\item Reidentification
			\item Encryption
			\item Hashing
		\end{enumerate}
	\end{frame}
	
	\begin{frame}{Q75: Explanation}
		\begin{block}{Correct Answer: B}
			When data is aggregated, semianonymous information or information that is not individually identifiable might become nonanonymous or identifiable in the process. This is reidentification. Privacy (reidentification) is a risk associated with big data analytics.
		\end{block}
		\begin{block}{Option Analysis}
			\begin{itemize}
				\item \textbf{A - Incorrect:} Data minimization reduces privacy risk
				\item \textbf{B - Correct:} Reidentification occurs through aggregation
				\item \textbf{C - Incorrect:} Encryption protects confidentiality
				\item \textbf{D - Incorrect:} Hashing protects integrity
			\end{itemize}
		\end{block}
	\end{frame}
	
	\begin{frame}{Q76: Big Data - Data Deluge}
		\begin{block}{Question}
			What challenge does "data deluge" present for big data initiatives?
		\end{block}
		\begin{enumerate}[(A)]
			\item It improves data quality
			\item It can result in quality problems, and effective decisions cannot be made based on poor quality data
			\item It reduces the need for governance
			\item It eliminates the need for skilled staff
		\end{enumerate}
	\end{frame}
	
	\begin{frame}{Q76: Explanation}
		\begin{block}{Correct Answer: B}
			Data deluge can result in quality problems. Big data is of no value if it cannot be harnessed for better business decision making. Effective decisions cannot be made based on poor quality data. In addition, more complex and unstructured data is being introduced, which can strain data governance efforts.
		\end{block}
		\begin{block}{Option Analysis}
			\begin{itemize}
				\item \textbf{A - Incorrect:} Data deluge often degrades quality
				\item \textbf{B - Correct:} Quality problems impede decision-making
				\item \textbf{C - Incorrect:} It increases governance needs
				\item \textbf{D - Incorrect:} It increases the need for skilled staff
			\end{itemize}
		\end{block}
	\end{frame}
	
	\begin{frame}{Q77: Blockchain Transparency}
		\begin{block}{Question}
			What blockchain characteristic facilitates extraordinary transparency in transaction logging?
		\end{block}
		\begin{enumerate}[(A)]
			\item Centralized processing
			\item The absence of centralization
			\item Symmetric encryption only
			\item Manual verification
		\end{enumerate}
	\end{frame}
	
	\begin{frame}{Q77: Explanation}
		\begin{block}{Correct Answer: B}
			The absence of centralization facilitates extraordinary transparency in transaction logging that is inherent to the blockchain structure. The chain structure makes it very difficult to falsify a blockchain entry because each entry contains both its own digital hash and the hash of the block before it.
		\end{block}
		\begin{block}{Option Analysis}
			\begin{itemize}
				\item \textbf{A - Incorrect:} Blockchain is decentralized
				\item \textbf{B - Correct:} Absence of centralization provides transparency
				\item \textbf{C - Incorrect:} Uses cryptographic hashing, not just symmetric
				\item \textbf{D - Incorrect:} Verification is automated through consensus
			\end{itemize}
		\end{block}
	\end{frame}
	
	\begin{frame}{Q78: Blockchain Potential Beyond Cryptocurrency}
		\begin{block}{Question}
			Which of the following is a potential application of blockchain technology beyond cryptocurrency?
		\end{block}
		\begin{enumerate}[(A)]
			\item Entertainment only
			\item Healthcare, property titles, and voting
			\item Only financial trading
			\item No other applications exist
		\end{enumerate}
	\end{frame}
	
	\begin{frame}{Q78: Explanation}
		\begin{block}{Correct Answer: B}
			Blockchain has been touted as a way to revolutionize banking, healthcare, property titles, and even voting. Some of the promises made may prove to be baseless hype. However, risk practitioners in industries affected by blockchain should take care to remain apprised of what developments are underway.
		\end{block}
		\begin{block}{Option Analysis}
			\begin{itemize}
				\item \textbf{A - Incorrect:} Blockchain has many applications beyond entertainment
				\item \textbf{B - Correct:} Healthcare, property titles, voting
				\item \textbf{C - Incorrect:} Blockchain has applications beyond finance
				\item \textbf{D - Incorrect:} Many applications exist
			\end{itemize}
		\end{block}
	\end{frame}
	
	\begin{frame}{Q79: AI Machine Learning Oversight}
		\begin{block}{Question}
			Why is oversight especially important for machine learning (ML) systems?
		\end{block}
		\begin{enumerate}[(A)]
			\item ML systems never change their behavior
			\item ML algorithms refine themselves over time based on results, and bad habits may be difficult to break
			\item ML systems are always perfectly accurate
			\item ML systems do not require human intervention
		\end{enumerate}
	\end{frame}
	
	\begin{frame}{Q79: Explanation}
		\begin{block}{Correct Answer: B}
			Oversight of AI is especially important in cases involving ML, algorithms designed to refine themselves over time based on results. Where "improvement" is poorly or incorrectly defined, iteration aimed at refining outcomes may result in lasting deviation from goals. Heuristic systems used in intrusion detection are one area of concern—an errant determination that malicious traffic is normal could result in willful blindness to an intrusion.
		\end{block}
		\begin{block}{Option Analysis}
			\begin{itemize}
				\item \textbf{A - Incorrect:} ML systems change their behavior
				\item \textbf{B - Correct:} Self-refining algorithms may deviate from goals
				\item \textbf{C - Incorrect:} ML systems can be inaccurate
				\item \textbf{D - Incorrect:} Human oversight is essential
			\end{itemize}
		\end{block}
	\end{frame}
	
	\begin{frame}{Q80: AI and Sensitive Data}
		\begin{block}{Question}
			What privacy and intellectual property risk may be overlooked when using AI?
		\end{block}
		\begin{enumerate}[(A)]
			\item AI never uses data for training
			\item The use of sensitive and personally identifiable information (PII) to train AI models
			\item AI only uses public data
			\item AI eliminates all privacy concerns
		\end{enumerate}
	\end{frame}
	
	\begin{frame}{Q80: Explanation}
		\begin{block}{Correct Answer: B}
			The use of sensitive and personally identifiable information (PII) to train AI models presents privacy and intellectual property risk that may be overlooked using standard risk identification and assessment processes. The risk practitioner should ensure that privacy impact assessments (PIAs) and fundamental rights impact assessments (FRIAs) are part of the AI risk assessment process.
		\end{block}
		\begin{block}{Option Analysis}
			\begin{itemize}
				\item \textbf{A - Incorrect:} AI requires large amounts of training data
				\item \textbf{B - Correct:} PII in training data creates risk
				\item \textbf{C - Incorrect:} AI can use proprietary data
				\item \textbf{D - Incorrect:} AI can increase privacy concerns
			\end{itemize}
		\end{block}
	\end{frame}
	
	% ================ SECTION 7: SECURITY FRAMEWORKS & ACCESS CONTROL ================
	
	\section{Security Frameworks \& Access Control}
	
	\begin{frame}{Q81: CIA Triad Trade-offs}
		\begin{block}{Question}
			Why is it impossible to simultaneously maximize all three principles of the CIA triad?
		\end{block}
		\begin{enumerate}[(A)]
			\item They are unrelated
			\item Typically there are trade-offs between the three principles, making it impossible to increase one without decreasing another or substantially increasing cost
			\item They are always perfectly aligned
			\item Only one principle matters
		\end{enumerate}
	\end{frame}
	
	\begin{frame}{Q81: Explanation}
		\begin{block}{Correct Answer: B}
			There are typically trade-offs between the three principles in any system, making it impossible to simultaneously increase one of them without either decreasing at least one of the others or substantially increasing cost. For instance, increasing confidentiality typically increases processing time, which reduces availability.
		\end{block}
		\begin{block}{Option Analysis}
			\begin{itemize}
				\item \textbf{A - Incorrect:} They are closely related
				\item \textbf{B - Correct:} Trade-offs are inherent in the CIA triad
				\item \textbf{C - Incorrect:} They often conflict
				\item \textbf{D - Incorrect:} All three principles matter
			\end{itemize}
		\end{block}
	\end{frame}
	
	\begin{frame}{Q82: Confidentiality - Need to Know}
		\begin{block}{Question}
			What principle ensures that individuals are given access only to information they need to perform their job functions?
		\end{block}
		\begin{enumerate}[(A)]
			\item Least privilege
			\item Need to know
			\item Availability
			\item Nonrepudiation
		\end{enumerate}
	\end{frame}
	
	\begin{frame}{Q82: Explanation}
		\begin{block}{Correct Answer: B}
			The need-to-know principle ensures that individuals are given access only to information that they need to perform their job functions. For example, most users do not need to see credit card numbers; therefore, the numbers are masked or hidden except for the last few digits. Least privilege refers to the minimum level of access needed, such as read-only versus edit.
		\end{block}
		\begin{block}{Option Analysis}
			\begin{itemize}
				\item \textbf{A - Incorrect:} Least privilege refers to the level of access
				\item \textbf{B - Correct:} Need to know limits access to needed information
				\item \textbf{C - Incorrect:} Availability is about access when needed
				\item \textbf{D - Incorrect:} Nonrepudiation is about proving origin
			\end{itemize}
		\end{block}
	\end{frame}
	
	\begin{frame}{Q83: Integrity - Protection Against Modification}
		\begin{block}{Question}
			Integrity refers to protection against:
		\end{block}
		\begin{enumerate}[(A)]
			\item Unauthorized disclosure
			\item Improper modification, exclusion, or destruction of information
			\item Unavailability of data when needed
			\item Denial of an action
		\end{enumerate}
	\end{frame}
	
	\begin{frame}{Q83: Explanation}
		\begin{block}{Correct Answer: B}
			Integrity refers to protection against improper modification, exclusion, or destruction of information and applies to actions taken by both authorized and unauthorized users as well as processes or activities operating on the system. Whenever data is changed in a manner other than that intended by the data owner, integrity is compromised.
		\end{block}
		\begin{block}{Option Analysis}
			\begin{itemize}
				\item \textbf{A - Incorrect:} Unauthorized disclosure is a confidentiality breach
				\item \textbf{B - Correct:} Improper modification is an integrity breach
				\item \textbf{C - Incorrect:} Unavailability is an availability breach
				\item \textbf{D - Incorrect:} Denial is a nonrepudiation issue
			\end{itemize}
		\end{block}
	\end{frame}
	
	\begin{frame}{Q84: Availability - 100\% Uptime}
		\begin{block}{Question}
			Why is an initial answer from process owners that systems require 100\% availability often inaccurate?
		\end{block}
		\begin{enumerate}[(A)]
			\item Because process owners always exaggerate
			\item Because the financial costs associated with guaranteeing true 100\% uptime are rapidly apparent
			\item Because availability is never important
			\item Because systems are always available
		\end{enumerate}
	\end{frame}
	
	\begin{frame}{Q84: Explanation}
		\begin{block}{Correct Answer: B}
			Risk practitioners should anticipate an initial answer from process owners that systems and data are essential and require 100\% availability. In most cases, this is not accurate, as rapidly becomes clear when the financial costs associated with guaranteeing true 100\% uptime are presented. From there, the risk practitioner should work with the business to determine specific requirements and align these with what the organization is willing to pay.
		\end{block}
		\begin{block}{Option Analysis}
			\begin{itemize}
				\item \textbf{A - Incorrect:} Process owners are not exaggerating intentionally
				\item \textbf{B - Correct:} Cost of 100\% uptime becomes apparent
				\item \textbf{C - Incorrect:} Availability is important
				\item \textbf{D - Incorrect:} Systems are not always available
			\end{itemize}
		\end{block}
	\end{frame}
	
	\begin{frame}{Q85: Nonrepudiation Controls}
		\begin{block}{Question}
			Which controls provide nonrepudiation?
		\end{block}
		\begin{enumerate}[(A)]
			\item Shared or generic logon credentials
			\item Digital signatures and certificate-based authentication in a PKI
			\item Simple usernames and passwords
			\item Antivirus software
		\end{enumerate}
	\end{frame}
	
	\begin{frame}{Q85: Explanation}
		\begin{block}{Correct Answer: B}
			Digital signatures and certificate-based authentication in a public key infrastructure (PKI) are examples of controls whose proper implementation and use can provide nonrepudiation. Shared or generic logon credentials provide no assurance. Individual usernames and passwords provide relatively weak assurance because these can be shared or guessed.
		\end{block}
		\begin{block}{Option Analysis}
			\begin{itemize}
				\item \textbf{A - Incorrect:} Shared credentials provide no nonrepudiation
				\item \textbf{B - Correct:} Digital signatures and PKI provide nonrepudiation
				\item \textbf{C - Incorrect:} Usernames/passwords provide weak assurance
				\item \textbf{D - Incorrect:} Antivirus does not provide nonrepudiation
			\end{itemize}
		\end{block}
	\end{frame}
	
	\begin{frame}{Q86: System Authorization Purpose}
		\begin{block}{Question}
			What is the purpose of the system authorization process (also called security assessment and authorization)?
		\end{block}
		\begin{enumerate}[(A)]
			\item To install new software
			\item To explicitly authorize a system to operate based on a thorough review of its security
			\item To eliminate all risk from a system
			\item To replace the need for testing
		\end{enumerate}
	\end{frame}
	
	\begin{frame}{Q86: Explanation}
		\begin{block}{Correct Answer: B}
			The system authorization process serves the purpose of explicitly authorizing a system to operate prior to it being allowed to do so. During assessment, an information system is thoroughly reviewed to determine the security of its design, development, testing, deployment, and operations. The culmination is a report to a suitably empowered senior manager recommending whether to authorize operation of the system.
		\end{block}
		\begin{block}{Option Analysis}
			\begin{itemize}
				\item \textbf{A - Incorrect:} Authorization is not about installing software
				\item \textbf{B - Correct:} Explicit authorization based on security review
				\item \textbf{C - Incorrect:} Risk cannot be eliminated entirely
				\item \textbf{D - Incorrect:} Testing is part of the assessment
			\end{itemize}
		\end{block}
	\end{frame}
	
	\begin{frame}{Q87: Authorizing Official Responsibility}
		\begin{block}{Question}
			Who is responsible for the formal decision to accept the risk and authorize a system for operational use?
		\end{block}
		\begin{enumerate}[(A)]
			\item The security evaluator
			\item The authorizing official (AO)
			\item The system administrator
			\item The end users
		\end{enumerate}
	\end{frame}
	
	\begin{frame}{Q87: Explanation}
		\begin{block}{Correct Answer: B}
			Authorization is the official decision by a senior manager—sometimes called an authorizing official (AO)—to approve an information system for operational use. This decision is the explicit acceptance of the risk documented in the evaluator's report. The decision to accept risk is ultimately their responsibility as a senior manager accountable for business matters, not the evaluator who performed the assessment.
		\end{block}
		\begin{block}{Option Analysis}
			\begin{itemize}
				\item \textbf{A - Incorrect:} The evaluator assesses, not authorizes
				\item \textbf{B - Correct:} The AO accepts risk and authorizes
				\item \textbf{C - Incorrect:} Administrators operate systems
				\item \textbf{D - Incorrect:} End users use the system
			\end{itemize}
		\end{block}
	\end{frame}
	
	\begin{frame}{Q88: Separation of Duties}
		\begin{block}{Question}
			What is the purpose of separation of duties (SoD)?
		\end{block}
		\begin{enumerate}[(A)]
			\item To improve system performance
			\item To prevent or detect errors and irregularities by assigning responsibility for initiating and recording transactions and custody of assets to separate individuals
			\item To increase the number of administrators
			\item To reduce IT costs
		\end{enumerate}
	\end{frame}
	
	\begin{frame}{Q88: Explanation}
		\begin{block}{Correct Answer: B}
			SoD is a basic internal control that prevents or detects errors and irregularities by assigning responsibility for initiating and recording transactions and the custody of assets to separate individuals. Most often, this takes the form of requiring one person to review and approve the work of another person.
		\end{block}
		\begin{block}{Option Analysis}
			\begin{itemize}
				\item \textbf{A - Incorrect:} SoD does not improve performance
				\item \textbf{B - Correct:} SoD prevents and detects errors
				\item \textbf{C - Incorrect:} SoD is not about increasing headcount
				\item \textbf{D - Incorrect:} SoD does not reduce costs
			\end{itemize}
		\end{block}
	\end{frame}
	
	\begin{frame}{Q89: SoD and Collusion}
		\begin{block}{Question}
			Why does separation of duties not guarantee security?
		\end{block}
		\begin{enumerate}[(A)]
			\item Because it is too expensive
			\item Because it can be circumvented when all participants in the process agree to work together to bypass the intended effect (collusion)
			\item Because it always fails
			\item Because it is not a real control
		\end{enumerate}
	\end{frame}
	
	\begin{frame}{Q89: Explanation}
		\begin{block}{Correct Answer: B}
			SoD does not guarantee security because it can be circumvented when all participants in the process agree to work together to bypass the intended effect of the control (collusion). However, because it is difficult to establish collusion and especially challenging to maintain it over time, effective SoD increases the odds that an incorrect or fraudulent transaction will be detected and addressed.
		\end{block}
		\begin{block}{Option Analysis}
			\begin{itemize}
				\item \textbf{A - Incorrect:} Cost is not the reason
				\item \textbf{B - Correct:} Collusion can circumvent SoD
				\item \textbf{C - Incorrect:} SoD is effective when implemented well
				\item \textbf{D - Incorrect:} SoD is a real and important control
			\end{itemize}
		\end{block}
	\end{frame}
	
	\begin{frame}{Q90: Job Rotation as Security Control}
		\begin{block}{Question}
			Why is job rotation an effective security mechanism?
		\end{block}
		\begin{enumerate}[(A)]
			\item It improves productivity
			\item It changes who is involved in which roles at which times, reducing the potential for collusion and increasing the odds that prior collusion can be detected
			\item It reduces training costs
			\item It eliminates all fraud
		\end{enumerate}
	\end{frame}
	
	\begin{frame}{Q90: Explanation}
		\begin{block}{Correct Answer: B}
			Job rotation is typically associated with business continuity and enterprise resiliency, but it is also an effective security mechanism because it changes who is involved in which roles at which times, reducing the potential for collusion and increasing the odds that prior collusion can be detected by someone who is not part of the scheme. In environments where there is substantial concern regarding collusion, job rotation may take the form of mandatory vacations.
		\end{block}
		\begin{block}{Option Analysis}
			\begin{itemize}
				\item \textbf{A - Incorrect:} Rotation can decrease efficiency during transitions
				\item \textbf{B - Correct:} Rotation reduces collusion potential
				\item \textbf{C - Incorrect:} Rotation increases training costs
				\item \textbf{D - Incorrect:} Rotation does not eliminate all fraud
			\end{itemize}
		\end{block}
	\end{frame}
	
	\begin{frame}{Q91: Identification Purpose}
		\begin{block}{Question}
			What is the purpose of identification in the access control process?
		\end{block}
		\begin{enumerate}[(A)]
			\item To prove the user's identity
			\item To assign a unique identifier to every individual, device, or process that has access to a system, enabling per-user tracking and logging
			\item To determine what resources a user can access
			\item To record user actions for auditing
		\end{enumerate}
	\end{frame}
	
	\begin{frame}{Q91: Explanation}
		\begin{block}{Correct Answer: B}
			With a unique identifier assigned to every individual, device, or process that has access to a system, it is possible to track and log activity on a per-user basis, creating the possibility of a targeted investigation were a problem to arise. Identification is often provided through a user ID, customer account number, employee identification number, or other unique element.
		\end{block}
		\begin{block}{Option Analysis}
			\begin{itemize}
				\item \textbf{A - Incorrect:} Authentication proves identity
				\item \textbf{B - Correct:} Identification enables per-user tracking
				\item \textbf{C - Incorrect:} Authorization determines access
				\item \textbf{D - Incorrect:} Accountability records actions
			\end{itemize}
		\end{block}
	\end{frame}
	
	\begin{frame}{Q92: Authentication Factors}
		\begin{block}{Question}
			Which of the following correctly lists the three factors of authentication?
		\end{block}
		\begin{enumerate}[(A)]
			\item Something you know, something you have, something you are
			\item Something you want, something you need, something you see
			\item Identification, authorization, accountability
			\item Password, username, token
		\end{enumerate}
	\end{frame}
	
	\begin{frame}{Q92: Explanation}
		\begin{block}{Correct Answer: A}
			Authentication is done using one or more of three factors: something you know (e.g., user ID, password), something you have (e.g., token, authenticator), and something you are (e.g., fingerprint, iris scan). Multifactor authentication (MFA) combines more than one factor during the login process to establish strong authentication.
		\end{block}
		\begin{block}{Option Analysis}
			\begin{itemize}
				\item \textbf{A - Correct:} Know, have, are
				\item \textbf{B - Incorrect:} Not valid authentication factors
				\item \textbf{C - Incorrect:} These are IAAA components
				\item \textbf{D - Incorrect:} Password and username are knowledge factors; token is possession
			\end{itemize}
		\end{block}
	\end{frame}
	
	\begin{frame}{Q93: Password Authentication Weakness}
		\begin{block}{Question}
			Why is authentication by knowledge (e.g., passwords) considered weak over time?
		\end{block}
		\begin{enumerate}[(A)]
			\item Because passwords are never used
			\item Because time is the enemy of knowledge authentication—any knowledge question can eventually be figured out, allowing one person to log in as another
			\item Because passwords are always too complex
			\item Because passwords cannot be changed
		\end{enumerate}
	\end{frame}
	
	\begin{frame}{Q93: Explanation}
		\begin{block}{Correct Answer: B}
			Time is the enemy of knowledge authentication, because any knowledge question can eventually be figured out, allowing one person to log in as another. Therefore, changing passwords on a regular basis is recommended even in high-complexity environments. Knowledge is often subject to replay attacks, even when what is asked is complex.
		\end{block}
		\begin{block}{Option Analysis}
			\begin{itemize}
				\item \textbf{A - Incorrect:} Passwords are widely used
				\item \textbf{B - Correct:} Knowledge can eventually be compromised
				\item \textbf{C - Incorrect:} Complex passwords are used for extended periods
				\item \textbf{D - Incorrect:} Passwords can and should be changed
			\end{itemize}
		\end{block}
	\end{frame}
	
	\begin{frame}{Q94: Possession Authentication Downside}
		\begin{block}{Question}
			What is a key downside of possession-based authentication (e.g., tokens, smart cards)?
		\end{block}
		\begin{enumerate}[(A)]
			\item It is too secure
			\item If the means of authentication is lost, stolen, or damaged, the legitimate user may be denied access without easy recourse
			\item It is too cheap
			\item It cannot be used with mobile phones
		\end{enumerate}
	\end{frame}
	
	\begin{frame}{Q94: Explanation}
		\begin{block}{Correct Answer: B}
			A downside of possession-based authentication is that should the means of authentication be lost, stolen, or damaged, the legitimate user may be denied access without easy recourse. Where availability is a primary concern, the risk practitioner should take care to examine what workarounds may be in place, keeping in mind that any workaround also represents a potential area of vulnerability.
		\end{block}
		\begin{block}{Option Analysis}
			\begin{itemize}
				\item \textbf{A - Incorrect:} No authentication is too secure
				\item \textbf{B - Correct:} Loss of token denies legitimate access
				\item \textbf{C - Incorrect:} Possession authentication can be expensive
				\item \textbf{D - Incorrect:} Mobile phones enable possession authentication
			\end{itemize}
		\end{block}
	\end{frame}
	
	\begin{frame}{Q95: Biometric Authentication Concern}
		\begin{block}{Question}
			What is a significant concern with biometric authentication?
		\end{block}
		\begin{enumerate}[(A)]
			\item It is always inaccurate
			\item Biometric data may be subject to laws and regulations governing privacy, with potentially severe penalties should the data be compromised
			\item It cannot be used for access control
			\item It is too easy to fake
		\end{enumerate}
	\end{frame}
	
	\begin{frame}{Q95: Explanation}
		\begin{block}{Correct Answer: B}
			Biometric data may be subject to laws and regulations governing privacy, with potentially severe penalties in some jurisdictions should the data be compromised. Also, while it is generally assumed that a person would not lose characteristics in ways that might deny access, this assumption does not hold under all conditions—soldiers or workers in industrial environments might experience physical injuries that affect their eyes or limbs.
		\end{block}
		\begin{block}{Option Analysis}
			\begin{itemize}
				\item \textbf{A - Incorrect:} Biometrics are generally accurate
				\item \textbf{B - Correct:} Privacy regulations create legal risk
				\item \textbf{C - Incorrect:} Biometrics are used for access control
				\item \textbf{D - Incorrect:} Biometrics are difficult to fake
			\end{itemize}
		\end{block}
	\end{frame}
	
	\begin{frame}{Q96: Multifactor Authentication Dependency}
		\begin{block}{Question}
			What risk do many MFA implementations introduce by relying on personal smartphones?
		\end{block}
		\begin{enumerate}[(A)]
			\item They are too secure
			\item Authentication becomes reliant on access to devices outside organizational control, which may be lost, damaged, or stolen
			\item They eliminate all fraud
			\item They do not require passwords
		\end{enumerate}
	\end{frame}
	
	\begin{frame}{Q96: Explanation}
		\begin{block}{Correct Answer: B}
			Many MFA implementations combine traditional passwords with codes sent to or generated by applications on personal smartphones. These arrangements, while convenient, make authentication reliant on access to devices outside organizational control. Should a mobile device required for MFA be unavailable, the result is likely to be a temporary inability of a valid user to authenticate.
		\end{block}
		\begin{block}{Option Analysis}
			\begin{itemize}
				\item \textbf{A - Incorrect:} No MFA is too secure
				\item \textbf{B - Correct:} Dependency on personal devices
				\item \textbf{C - Incorrect:} MFA does not eliminate fraud
				\item \textbf{D - Incorrect:} MFA often still requires passwords
			\end{itemize}
		\end{block}
	\end{frame}
	
	\begin{frame}{Q97: Authorization - Least Privilege}
		\begin{block}{Question}
			What does proper employment of least privilege ensure in authorization?
		\end{block}
		\begin{enumerate}[(A)]
			\item That all users have full control over data
			\item That employees only have the level of access they require to perform their job functions
			\item That access is never restricted
			\item That all users have administrative access
		\end{enumerate}
	\end{frame}
	
	\begin{frame}{Q97: Explanation}
		\begin{block}{Correct Answer: B}
			Authorization refers to the privileges or permissions the person will have. Proper employment of least privilege is important to ensure that employees only have the level of access they require to perform their job functions. Permissions may be applied on a per-file or per-folder basis.
		\end{block}
		\begin{block}{Option Analysis}
			\begin{itemize}
				\item \textbf{A - Incorrect:} Full control violates least privilege
				\item \textbf{B - Correct:} Least privilege limits access to what is needed
				\item \textbf{C - Incorrect:} Access must be controlled
				\item \textbf{D - Incorrect:} Not all users need administrative access
			\end{itemize}
		\end{block}
	\end{frame}
	
	\begin{frame}{Q98: Authorization Isolation}
		\begin{block}{Question}
			How does authorization isolation based on time or location improve security?
		\end{block}
		\begin{enumerate}[(A)]
			\item It allows access from anywhere at any time
			\item It prevents a user from using account privileges at times or locations when they should not be required, reducing the potential scope of monitoring
			\item It eliminates the need for authentication
			\item It gives users more privileges
		\end{enumerate}
	\end{frame}
	
	\begin{frame}{Q98: Explanation}
		\begin{block}{Correct Answer: B}
			Authorization is typically granted only for a limited period. This isolation on a time- or location-delineated basis prevents a user from using account privileges at times when they should not be required and allows an organization to maximize the effectiveness of its controls by reducing the potential scope of monitoring. For example, access might only be allowed during normal working hours.
		\end{block}
		\begin{block}{Option Analysis}
			\begin{itemize}
				\item \textbf{A - Incorrect:} Isolation restricts access
				\item \textbf{B - Correct:} Time/location isolation reduces misuse potential
				\item \textbf{C - Incorrect:} Authentication is still required
				\item \textbf{D - Incorrect:} Isolation restricts, not expands, privileges
			\end{itemize}
		\end{block}
	\end{frame}
	
	\begin{frame}{Q99: Accountability - Log Protection}
		\begin{block}{Question}
			Why should audit logs not be changeable by operations staff, even administrators?
		\end{block}
		\begin{enumerate}[(A)]
			\item To save storage space
			\item To prevent a user from disabling the logging function or altering log data that implicates their own account in suspicious or malicious activity
			\item To improve system performance
			\item Because logs contain no useful information
		\end{enumerate}
	\end{frame}
	
	\begin{frame}{Q99: Explanation}
		\begin{block}{Correct Answer: B}
			To prevent a user from disabling the logging function or altering log data that implicates their own account in suspicious or malicious activity, audit logs should not be changeable by operations staff—even administrators. Logs also may contain sensitive data, so even view access to log content should be strictly controlled based on clear business need.
		\end{block}
		\begin{block}{Option Analysis}
			\begin{itemize}
				\item \textbf{A - Incorrect:} Storage is not the primary concern
				\item \textbf{B - Correct:} Prevents tampering with evidence
				\item \textbf{C - Incorrect:} Performance is not the reason
				\item \textbf{D - Incorrect:} Logs contain valuable evidence
			\end{itemize}
		\end{block}
	\end{frame}
	
	\begin{frame}{Q100: Symmetric Key Distribution}
		\begin{block}{Question}
			What is the primary challenge with symmetric key cryptography?
		\end{block}
		\begin{enumerate}[(A)]
			\item It is too slow for bulk data
			\item The need for one party to deliver keys to another with whom it would like to exchange data, which is difficult in environments like e-commerce with unknown parties
			\item It uses too much processing power
			\item It cannot encrypt large messages
		\end{enumerate}
	\end{frame}
	
	\begin{frame}{Q100: Explanation}
		\begin{block}{Correct Answer: B}
			Symmetric key cryptography has two major disadvantages. One is the need for one party to deliver keys to another with whom it would like to exchange data. Where communications are established between trusted parties, this may be workable through physical methods. However, in environments such as e-commerce, customers are unknown and not trusted entities, making brokering a connection more difficult.
		\end{block}
		\begin{block}{Option Analysis}
			\begin{itemize}
				\item \textbf{A - Incorrect:} Symmetric is fast for bulk data
				\item \textbf{B - Correct:} Key distribution is the main challenge
				\item \textbf{C - Incorrect:} Symmetric uses less processing power than asymmetric
				\item \textbf{D - Incorrect:} Symmetric encrypts large messages efficiently
			\end{itemize}
		\end{block}
	\end{frame}
	
	\begin{frame}{Q101: Asymmetric Key Relationship}
		\begin{block}{Question}
			In asymmetric encryption, what is the relationship between the public and private keys?
		\end{block}
		\begin{enumerate}[(A)]
			\item They are identical
			\item They are mathematically related, but it is computationally infeasible to determine one from the other
			\item They have no relationship
			\item The public key can decrypt what the private key encrypts, but not vice versa
		\end{enumerate}
	\end{frame}
	
	\begin{frame}{Q101: Explanation}
		\begin{block}{Correct Answer: B}
			In a Diffie-Hellman setup, two mathematically related keys are created: a private key and a public key. It is computationally infeasible to determine the value of one key from the other, but the manner of their generation means that each has the effect, when used in an appropriate asymmetric algorithm, of creating ciphertext that only the other has the ability to return to plaintext.
		\end{block}
		\begin{block}{Option Analysis}
			\begin{itemize}
				\item \textbf{A - Incorrect:} Keys are different
				\item \textbf{B - Correct:} Mathematically related but infeasible to derive
				\item \textbf{C - Incorrect:} They are mathematically related
				\item \textbf{D - Incorrect:} Each can decrypt what the other encrypts
			\end{itemize}
		\end{block}
	\end{frame}
	
	\begin{frame}{Q102: Asymmetric Algorithm Usage}
		\begin{block}{Question}
			Why are asymmetric algorithms typically used only to encrypt short messages?
		\end{block}
		\begin{enumerate}[(A)]
			\item They are less secure than symmetric algorithms
			\item They are computationally intensive and slow relative to symmetric algorithms
			\item They cannot encrypt any data
			\item They require no keys
		\end{enumerate}
	\end{frame}
	
	\begin{frame}{Q102: Explanation}
		\begin{block}{Correct Answer: B}
			Asymmetric algorithms are computationally intensive and slow relative to symmetric algorithms. For that reason, asymmetric cryptography is typically used only to encrypt short messages. The most common use of asymmetric algorithms is to distribute symmetric keys that can then be used by the participants for fast, secure communication.
		\end{block}
		\begin{block}{Option Analysis}
			\begin{itemize}
				\item \textbf{A - Incorrect:} Asymmetric is not less secure
				\item \textbf{B - Correct:} Asymmetric is slow and intensive
				\item \textbf{C - Incorrect:} Asymmetric can encrypt data
				\item \textbf{D - Incorrect:} Asymmetric requires key pairs
			\end{itemize}
		\end{block}
	\end{frame}
	
	\begin{frame}{Q103: Hashing Purpose}
		\begin{block}{Question}
			What is the original purpose of hashing algorithms?
		\end{block}
		\begin{enumerate}[(A)]
			\item To encrypt messages for confidentiality
			\item To detect accidental errors caused by interference or network problems
			\item To provide nonrepudiation
			\item To replace digital signatures
		\end{enumerate}
	\end{frame}
	
	\begin{frame}{Q103: Explanation}
		\begin{block}{Correct Answer: B}
			Hashing by itself provides effective protection against accidental changes in a message, such as those caused by interference, and this was its original purpose. When a sender wants to ensure a message is not affected by noise or network problems, he or she can hash the message and send the digest along with the message. The receiver can then hash the received message and verify that the resulting digest matches.
		\end{block}
		\begin{block}{Option Analysis}
			\begin{itemize}
				\item \textbf{A - Incorrect:} Hashing does not provide confidentiality
				\item \textbf{B - Correct:} Hashing detects accidental errors
				\item \textbf{C - Incorrect:} Hashing alone does not provide nonrepudiation
				\item \textbf{D - Incorrect:} Digital signatures use hashing, not replace it
			\end{itemize}
		\end{block}
	\end{frame}
	
	\begin{frame}{Q104: Hashing Limitation}
		\begin{block}{Question}
			Why is hashing alone insufficient to protect against intentional changes to a message?
		\end{block}
		\begin{enumerate}[(A)]
			\item Because hashes are too long
			\item Because a malicious user who intercepts a transmission can create a new digest after altering the message
			\item Because hashes cannot be calculated
			\item Because hashes are always wrong
		\end{enumerate}
	\end{frame}
	
	\begin{frame}{Q104: Explanation}
		\begin{block}{Correct Answer: B}
			A malicious user who intercepts a transmission that includes a hash will not be deterred by hashing, because they can create a new digest after altering the message. To guarantee integrity against intentional and accidental changes, cryptologists extended the basic idea of hashing to include features of public key encryption, resulting in digital signatures.
		\end{block}
		\begin{block}{Option Analysis}
			\begin{itemize}
				\item \textbf{A - Incorrect:} Hash length is not the issue
				\item \textbf{B - Correct:} Attacker can recompute the hash
				\item \textbf{C - Incorrect:} Hashes are easily calculated
				\item \textbf{D - Incorrect:} Hashes are reliable when computed correctly
			\end{itemize}
		\end{block}
	\end{frame}
	
	\begin{frame}{Q105: Digital Signature Nonrepudiation}
		\begin{block}{Question}
			Why does a digital signature provide nonrepudiation?
		\end{block}
		\begin{enumerate}[(A)]
			\item Because anyone can create the signature
			\item Because only the sender's private key could have encrypted the digest, so the sender cannot plausibly deny sending it
			\item Because the signature is symmetric
			\item Because the signature is public
		\end{enumerate}
	\end{frame}
	
	\begin{frame}{Q105: Explanation}
		\begin{block}{Correct Answer: B}
			A digital signature combines a hash function with the ability of public key encryption to prove the author's identity. The recipient's ability to decrypt the digital signature with the sender's public key means that it must have been encrypted with the sender's private key and is, therefore, the genuine digest for the message. Moreover, the sender cannot plausibly deny it, because only the sender's private key could have encrypted it.
		\end{block}
		\begin{block}{Option Analysis}
			\begin{itemize}
				\item \textbf{A - Incorrect:} Only the private key holder can sign
				\item \textbf{B - Correct:} Private key encryption proves origin
				\item \textbf{C - Incorrect:} Digital signatures use asymmetric encryption
				\item \textbf{D - Incorrect:} The signature is not simply public
			\end{itemize}
		\end{block}
	\end{frame}
	
	% ================ FINAL SLIDES ================
	
	\begin{frame}{Security Concepts Key Takeaways}
		\begin{block}{Essential Concepts}
			\begin{itemize}
				\item CIA triad involves trade-offs; you can't maximize all three
				\item Need-to-know limits access to needed information; least privilege limits access level
				\item Nonrepudiation requires digital signatures or PKI certificates
				\item SoD prevents errors but can be circumvented by collusion
				\item Job rotation reduces collusion potential
				\item Authentication factors: know, have, are
				\item Knowledge authentication weakens over time; possession has availability downsides
				\item Biometrics carry privacy regulation risk
				\item Symmetric key distribution is the main challenge; asymmetric is slow
				\item Hashing protects integrity but cannot prove origin; digital signatures provide nonrepudiation
			\end{itemize}
		\end{block}
	\end{frame}
	
	\begin{frame}{Final Advice}
		\begin{block}{Exam Preparation}
			\centering
			\vspace{0.3cm}
			\textbf{Remember:}
			\begin{itemize}
				\item Understand the \textit{why} behind each technology risk concept
				\item Know the difference between similar terms
				\item Understand how emerging technologies introduce new risk
				\item Practice scenario-based questions
				\item Use the ISACA mindset
			\end{itemize}
			\vspace{0.5cm}
			\textbf{Good Luck on Your CRISC Exam!}
		\end{block}
	\end{frame}
	
\end{document}